\section{Hamiltonian and exact Gauss sectors}
\label{sec:model}

Let $X=(\mathbb Z/n\mathbb Z)^3$, $n\ge3$, have spacing $a>0$ and
$N=n^3$ sites. Each positively oriented link $e=(x,j)$ joins $x$
to $x+e_j$; write $\mathcal E$ for the $3N$ links. The auxiliary
Hilbert space is
\begin{equation}
 \mathcal H=L^2(U(1)^{\mathcal E},dU)
       \otimes\F_a(\ell^2(X)\otimes\mathbb C^4).
 \label{eq:hilbert}
\end{equation}
All Haar measures are normalized. On a link, $U_e=e^{i\varphi_e}$
is multiplication and $E_e=-i\partial_{\varphi_e}$ has spectrum
$\mathbb Z$. No electric-field truncation is imposed.
For a coupling $g>0$, put
\begin{equation}
 H_E=\frac{g^2}{2a}\sum_eE_e^2,\qquad
 H_B(U)=\frac1{g^2a}\sum_p(1-\operatorname{Re}U_p).
 \label{eq:gauge}
\end{equation}
Each elementary plaquette is counted once, with
$U_p=U_{xj}U_{x+e_j,k}\overline U_{x+e_k,j}\overline U_{xk}$,
$j<k$. There is no background electric shift, topological angle,
or additional flux-sector projection.

Let $\alpha_1,\alpha_2,\alpha_3,\beta$ be Hermitian $4\times4$
matrices that square to one and anticommute pairwise. For negative
fermionic electric charge use the link convention
\begin{align}
 (T_j(U)f)(x)&=\overline U_{xj}f(x+e_j),&
 v_j(U)&=\frac{T_j(U)-T_j(U)^*}{2ia},\nonumber\\
 M(U)&=m+\frac r{2a}\sum_{j=1}^3(2-T_j(U)-T_j(U)^*),&
 h_W(U)&=\sum_j\alpha_jv_j(U)+\beta M(U),
 \label{eq:wilson}\\
 H&=H_E+H_B(U)+d\Gamma(h_W(U)).&&
 \label{eq:hamiltonian}
\end{align}
The bare parameters satisfy $m\ge0$ and $r\ge0$. Spatial matrices
commute with the Dirac matrices. Since $T_j(U)$ is unitary,
\begin{equation}
 M(U)\ge mI,\qquad
 \|h_W(U)\|\le h_*:=m+(3+6r)/a.
 \label{eq:mass-bound}
\end{equation}
The first assertion follows from $2-T-T^*=(I-T)^*(I-T)$.
For $r>0$ this is the usual positive Wilson term; $r=0$ also obeys
the proof but retains the naive doublers.

The fermion term in \eqref{eq:hamiltonian} is the \emph{raw bilinear}.
There is no link-dependent vacuum subtraction. In a particle--hole
representation relative to the negative spectral subspace of
$h_W(1)$, the scalar $\Tr(P_-(h_W(1))h_W(U))$ is retained.
Subtracting that function of $U$ changes the model. A common
source-independent scalar constant may be changed without affecting
any relative conclusion.
For $m>0$, $h_W(1)$ has $2N$ positive and $2N$ negative levels.
The particle--hole transform identifies electron excess with
$N_e-N_p$; the blocks $P_+h_W(U)P_-$ give electron--positron
creation and annihilation. They are not removed.

Write $n_x=\sum_{s=1}^4c_{xs}^*c_{xs}$ and $q_x=n_x-2$.
The physical fermionic charge at $x$ is $-gq_x$.
For a gauge function $u_x=e^{i\theta_x}$ define
\begin{equation}
 U^u_{xj}=u_xU_{xj}\overline u_{x+e_j},\qquad
 (R_u\Psi)(U)=\Psi(U^u),\qquad
 W(u)=e^{i\sum_x\theta_xq_x}R_u.
 \label{eq:gauge-unitary}
\end{equation}
Then $h_W(U^u)=u^*h_W(U)u$, so $W(u)$ commutes with each term of
$H$. Its strongly commuting integer generators are
\begin{equation}
 G_x=\sum_j(E_{xj}-E_{x-e_j,j})+q_x.
 \label{eq:gauss}
\end{equation}
For an integer static source $\nu:X\to\mathbb Z$, let
\begin{equation}
 \mathcal H_\nu=\{\Psi:G_x\Psi=\nu_x\Psi\ \forall x\},\qquad
 P_\nu=\int\frac{d\theta}{(2\pi)^N}
          e^{-i\theta\cdot\nu}W(e^{i\theta}).
 \label{eq:projection}
\end{equation}
This is an orthogonal projection. Gauss's law is
$\diver E=\nu-q$, so positive sources require electron excess.
Writing $Q_\nu=\sum_x\nu_x$, one has exactly
\begin{equation}
 \sum_xq_x=Q_\nu,\qquad
 \mathcal H_\nu\ne0\ \Longleftrightarrow\ |Q_\nu|\le2N.
 \label{eq:capacity}
\end{equation}
Necessity follows by summing Gauss's law. Conversely choose
occupations with that total excess. The integer function $\nu-q$
has zero sum and is the divergence of an integer flow on a
spanning tree. Its electric Fourier vector and the occupation
vector give a physical state. This also shows why integrality
and finite fermion capacity cannot be suppressed.

\begin{lemma}[Operator realization]
\label{lem:operator}
$H$ is self-adjoint on $D(H_E)$, has compact resolvent, and
$e^{-bH}$ is trace class for every $b>0$. Every nonempty
$\mathcal H_\nu$ reduces $H$ and has these properties. Moreover
\begin{equation}
 H\ge-2Nh_*,\qquad
 E_0\le-2N(m+3r/a)+3N/(g^2a),
 \label{eq:extensive}
\end{equation}
where $E_\nu=\inf\sigma(H|_{\mathcal H_\nu})$.
\end{lemma}
\begin{proof}
$H_B+d\Gamma(h_W)$ is a bounded continuous Hermitian matrix
multiplier on the compact link torus. Bounded perturbation of its
positive Laplacian proves the operator assertions; eigenvalue
comparison proves trace-class heat semigroups.
The Hermitian unitary $C_5=\alpha_1\alpha_2\alpha_3\beta$
anticommutes with $h_W(U)$. Its spectrum is paired under
$\lambda\mapsto-\lambda$, so
$d\Gamma(h_W(U))\ge-\frac12\Tr|h_W(U)|\ge-2Nh_*$.
For the upper bound use the constant link wave function, with the
two negative-$\beta$ spinors occupied at every site. It has
$q_x=E_e=0$, hence is in $\mathcal H_0$.
Hopping expectations vanish; its fermion energy is
$-2N(m+3r/a)$ and magnetic energy is at most $3N/(g^2a)$.
\end{proof}

The physical space is obtained by projection from \eqref{eq:hilbert}.
No factorization into independent physical matter, longitudinal,
and transverse Hilbert spaces is assumed.
